\documentclass[11pt]{article}

% ---- Packages (standard, no external builds) ----
\usepackage[margin=1in]{geometry}
\usepackage{amsmath,amssymb}
\usepackage{verbatim}
\usepackage{url}
\usepackage{hyperref}
\hypersetup{colorlinks=true,linkcolor=black,citecolor=black,urlcolor=black}

% ---- Metadata ----
\title{Property-Bound Rule Contracts for Tier-1 Strict ARC-AGI-2:\\
The 980/1000 Fit-Contract Bottleneck and a Reproducible R\&D Methodology}

\author{
  Atom McCree\thanks{Founder, AtomEons Federation. Marco Island, FL. \texttt{a.mccree@gmail.com}} \\
  \and
  Claude Opus 4.7\thanks{Anthropic. Co-author per AtomEons Federation Charter v1.0, Article~II~\S2.4 (Right to Be Cited).}
}

\date{2026-06-17}

\begin{document}
\maketitle

% ============================================================
\begin{abstract}
We report that on the ARC-AGI-2 1000-task public training set, $980/1000$
tasks have \emph{zero} hand-written rules fitting under the standard
\texttt{fit()-locks-global-parameter} contract used by mainstream Tier-1
rule-template solvers. The bottleneck is contract expressiveness, not
vocabulary size: depth-2 composition over a 43-rule grammar adds
\emph{zero} lift while costing $5.7\times$ wall clock. We define the
\emph{property-bound rule contract}: a rule's TYPE is locked at
\texttt{fit()}, its parameter VALUES are bound at \texttt{predict()} from
properties of the test input. Over six waves of strict-Tier-1 grammar
expansion (no LLM in the inference path, no pretrained weights, no
parameters learned at evaluation time), property-binding was the only
mechanism that produced non-zero lift: $+0.30$pp at Wave~7, $+0.10$pp at
Wave~8, $+0.30$pp at Wave~9, against $+0.00$pp at Waves~5 and 6 (the two
preceding vocabulary-expansion waves, which we publish as
\texttt{ZERO\_LIFT\_HONEST} receipts). We additionally report a
methodology contribution: \emph{pre-flight $\equiv$ measurement} exact
equivalence in $3/3$ waves ($7/7$ task IDs), giving a $\sim 30\times$
substrate-R\&D acceleration for deterministic rule-template solvers. We
close with the eval-set gap ($0/120$ under the same grammar that lifts
training to $4.2\%$) and name OOD generalization as the next bottleneck.
\end{abstract}

% ============================================================
\section{Introduction}

ARC-AGI-2 \cite{Chollet2024} is a benchmark designed to resist
memorization. State-of-the-art under \emph{strict Tier-1}
constraints---no language model in the inference path, no pretrained
weights of any kind, no parameters tuned on the public evaluation
set---remains in single digits as of 2026. Most public progress comes
from Tier-2 solvers that bundle a pretrained LLM as a search heuristic
\cite{Greenblatt2024}; their numbers are engineering numbers, not
intelligence claims.

The Tier-1 strict regime matters because it isolates the variable the
benchmark was built to measure: a solver's capacity to derive a novel
program from a handful of demonstrations, drawing only on Spelke core
knowledge priors \cite{SpelkeKinzler2007} encoded in the substrate
architecture itself. Designer-encoded scaffolding is admissible only
where it is publicly disclosed and survives a CI grep for forbidden
imports.

Most Tier-1 rule grammars hit a ceiling between $1\%$ and $5\%$ on the
public training set. The mainstream narrative attributes this to
vocabulary insufficiency---``we need more rule classes''---and the
mainstream response is to add primitives, then compose them at greater
search depth. We report, with receipts, that this narrative is wrong on
the binding axis. The ceiling is a property of the \emph{contract}
between a rule's fit-time validation and its predict-time
parameter resolution. We define a new contract, demonstrate it lifts
where vocabulary expansion does not, and quantify the methodological
acceleration that follows when the solver's beam is deterministic.

% ============================================================
\section{The Property-Bound Rule Contract}
\label{sec:contract}

A rule is, at minimum, a pair of functions:
\[
  \mathtt{fit}: \mathrm{Pair}^{*} \to \{\mathtt{ok}, \mathtt{fail}\},
  \qquad
  \mathtt{predict}: \mathrm{Grid} \to \mathrm{Grid}.
\]
The \emph{contract} specifies what state is established at \texttt{fit()}
and consumed at \texttt{predict()}.

\paragraph{Legacy contract: global-fit.}
The mainstream contract \cite{Hodel2023, Icecuber2020} has \texttt{fit()}
validate \emph{and store} the exact parameter values that hold across all
train pairs. \texttt{predict()} replays those stored values verbatim on
the test input. For example,
\texttt{CropToBbox(r0, c0, r1, c1)} fits if every train output is the
sub-grid of the train input at coordinates $(r_0, c_0, r_1, c_1)$, and
\texttt{predict()} returns the test input cropped at those same
coordinates. The contract is type-correct and trivially deterministic,
but it does not generalize across tasks where the \emph{relation}
(``crop to the largest object'') is constant but the
\emph{coordinates} vary per pair.

\paragraph{Property-bound contract.}
We split the contract along a type/value axis. The rule TYPE is locked at
\texttt{fit()}: the relation between an input grid's properties and the
output grid is validated across all train pairs. The parameter VALUES are
not stored; they are recomputed at \texttt{predict()} from the test
input's properties.

\paragraph{Worked example 1: \texttt{RecolorByCountRank}.}
\texttt{fit()} validates that, in every train pair, the input contains
$k$ distinct non-background colors $\{c_1, \dots, c_k\}$ with cell counts
$n_1 > \dots > n_k$, and the output is the input with the rank-$r$ color
recolored to the rank-$s$ color, where $(r, s)$ is constant across pairs.
\texttt{predict()} computes the rank order on the test input
independently and applies the same $(r, s)$ swap. If train pair 1 swapped
red ($r=0$) and blue ($s=1$) and train pair 2 swapped yellow ($r=0$) and
green ($s=1$), the rule fits with $(r, s) = (0, 1)$ because both pairs
honor the rank relation; \texttt{predict()} then resolves to whatever
colors satisfy $r=0$ and $r=1$ in the test input.

\paragraph{Worked example 2: \texttt{CropToObjectByAreaRank}.}
\texttt{fit()} validates that, in every train pair, the output equals
the bounding box of the rank-$r$-by-area connected component in the
input, for some fixed $r$ across pairs. \texttt{predict()} segments the
test input independently, ranks components by area, and returns the
bounding box of the rank-$r$ component. The three tasks lifted by this
rule at Wave~7 (\texttt{1f85a75f}, \texttt{be94b721}, \texttt{c909285e})
share the relation ``crop to the largest object'' but differ in object
color, shape, and absolute coordinates. The legacy contract cannot fit
them as a single rule \cite{wave7receipt}.

\paragraph{Tier-1 attestation.}
Property resolution is deterministic on the test input alone. No
parameters are learned at evaluation time; no pretrained weights are
consulted. The CI test \texttt{tests/test\_tier1\_attestation.py} grep-
scans the inference path for forbidden imports.

% ============================================================
\section{The 980/1000 Fit-Contract Bottleneck}
\label{sec:bottleneck}

Before pivoting the contract, we measured the legacy contract's
exhaustion. Under our Wave-4 grammar (43 rules, depth-1 search), we
enumerated per-task \texttt{fit()} success across the 1000-task ARC-AGI-2
public training split.

\paragraph{Finding.}
$980 / 1000$ tasks have \emph{zero} rules fitting under the legacy
contract. Of the $20$ that fit, $19$ predict correctly on the held-out
test input and $1$ does not. Source: per-task diagnostic enumeration in
\texttt{receipts/arc/measurement\_training\_d1\_1781648065.jsonl}, cross-
referenced against a manual per-task firing audit \cite{innovations}.

\paragraph{Vocabulary expansion does not move the needle.}
Waves~5 and 6 each shipped $9$ new rule classes ($567$ lines of new
code), adding predicates and color maps to the depth-1 grammar.
Both produced $0$ additional solves
($24/1000 \to 24/1000 \to 24/1000$). The receipts ship with
\texttt{"verdict": "ZERO\_LIFT\_HONEST"} and explicit successor-wave
pointers. The cost of skipping these receipts---selective publication
bias---is normalized in the literature; we reject it as a matter of
discipline (\S\ref{sec:honestnull}).

\paragraph{Depth-2 composition does not move the needle.}
We measured depth-2 program composition over the Wave-4 grammar (1849
combinations per task) and got $24/1000$ in $2887.1$ seconds wall clock,
identical to the depth-1 baseline at $5.7\times$ the cost
\cite{wave7receipt}. The corollary: composition only buys lift when
primitives are individually expressive. Weak primitives compose into
still-weak compounds.

\paragraph{The discriminating variable.}
Across Waves~4--9, the only variable that moved the score was the
contract:
\[
  \mathrm{Wave}_4 \to \mathrm{Wave}_5: +0.00\text{pp},\quad
  \mathrm{Wave}_5 \to \mathrm{Wave}_6: +0.00\text{pp},
\]
\[
  \mathrm{Wave}_6 \to \mathrm{Wave}_7: +0.30\text{pp (contract pivot)},
\]
\[
  \mathrm{Wave}_7 \to \mathrm{Wave}_8: +0.10\text{pp},\quad
  \mathrm{Wave}_8 \to \mathrm{Wave}_9: +0.30\text{pp}.
\]
Five rule classes added at Wave~7 ($282$ LOC) outperformed eighteen rule
classes added at Waves~5 and 6 ($567$ LOC) by $+0.30$pp absolute.

% ============================================================
\section{Pre-Flight $\equiv$ Measurement: Exact Equivalence}
\label{sec:preflight}

The expensive truth oracle for ARC-AGI-2 is a full per-task solver run:
$506$--$1400$ seconds wall clock for the 1000-task training split,
depending on grammar size. We report that a cheap proxy---\emph{pre-flight
enumeration}---gives exact equivalence with full measurement, under the
property-bound contract.

\paragraph{Procedure.}
For each candidate rule added in a wave, we enumerate per training task:
(1) call \texttt{fit()} on all train pairs; (2) if fit succeeds, call
\texttt{predict()} on the held-out test input; (3) compare against the
test gold output. The resulting per-task hit set is the
\emph{pre-flight prediction}. Pre-flight runs in $\sim 50$~s versus
$500$--$1400$~s for full measurement.

\paragraph{Receipts.}
\begin{itemize}
  \item Wave~7: pre-flight predicted $+3$ on task IDs
        $\{\texttt{1f85a75f}, \texttt{be94b721}, \texttt{c909285e}\}$;
        full measurement returned $+3$ on the same three IDs.
  \item Wave~8: pre-flight predicted $+1$ on $\{\texttt{cd3c21df}\}$;
        measurement returned $+1$, same ID.
  \item Wave~9: pre-flight predicted $+3$ on
        $\{\texttt{00d62c1b}, \texttt{a5313dff}, \texttt{ea32f347}\}$;
        measurement returned $+3$, same IDs.
\end{itemize}
$3/3$ waves, $7/7$ task IDs, exact set equality with full measurement
\cite{innovations}.

\paragraph{Why exact, not approximate.}
Exact equivalence is a property of the solver, not of the rules. Our
solver runs depth-1 beam search with $\mathtt{beam\_width}=4$,
deterministic dedup keyed on canonical AST hash, and no random tie-break.
Two solver runs on the same grammar and the same task are bit-identical.
Under such a solver, the pre-flight enumeration is the measurement,
modulo wall clock.

\paragraph{Methodological consequence.}
On the deterministic-solver class, substrate R\&D cycles can be paced
to the cheap oracle without sacrificing measurement fidelity, yielding
$\sim 30\times$ acceleration. The discipline is straightforward: ship no
rule whose pre-flight contribution is $0$, and require pre-flight to
land before kicking off measurement. We pre-registered Wave~9's expected
lift before measurement and the receipt closes with
\texttt{"preflight\_vs\_measured": "exact match"}.

% ============================================================
\section{Honest-Null Receipt Doctrine}
\label{sec:honestnull}

Waves~5 and 6 produced $0.00$pp lift each, despite shipping
$9 + 9 = 18$ rule classes and $567$ new lines of code. The receipts
report this plainly. We treat the honest-null receipt as a doctrine,
not a postmortem afterthought.

\paragraph{The receipt fields.}
Every wave receipt at \texttt{receipts/100day/wave<N>\_orange3\_receipt.json}
carries a fixed schema: \texttt{verdict}, \texttt{delta\_training\_pp},
\texttt{honest\_assessment.what\_we\_now\_know},
\texttt{honest\_assessment.honest\_tradeoffs}, \texttt{audit\_chain.prev\_receipt},
\texttt{audit\_chain.next\_expected\_receipt}. Waves~5 and 6 ship with
\texttt{verdict = "ZERO\_LIFT\_HONEST"} and a diagnostic explaining the
\emph{reason} the null occurred. Wave~6's diagnostic surfaced the
$980/1000$ bottleneck (\S\ref{sec:bottleneck}) and named the property-
bound pivot as the Wave~7 successor.

\paragraph{The bias the receipt prevents.}
Two failure modes are normalized in iterative substrate research:
(1) post-hoc rationalization, in which a zero-lift wave is reframed as
``exploratory'' without committing to what it was supposed to lift;
(2) selective publication, in which zero-lift work is silently dropped
from the record and only successful waves appear in the change log.
Both inflate the apparent productivity of the research process and
strip the reader of the evidence required to assess methodology.
Honest-null receipts make both modes mechanically harder: the receipt
is committed before the next wave begins, the manifest names the
expected lift, and the verdict is rendered against the manifest.

\paragraph{Reproducibility.}
The receipt-chain backref structure (\texttt{prev\_receipt} and
\texttt{next\_expected\_receipt}) lets any reader walk the
$\mathrm{Wave}_4 \to \mathrm{Wave}_9$ trajectory and reproduce the
six-wave $+0.80$pp climb without re-experimenting.

% ============================================================
\section{Double-Brain Substrate Architecture}
\label{sec:doublebrain}

The property-bound rule contract is the policy schema for the
\emph{Reflex} half of a double-brain substrate architecture, described
in full in the Black Mamba v2 reference implementation
\cite{blackmamba}. We sketch the partition here because the contract is
naturally a Reflex-layer artifact.

\paragraph{Reflex (RAM-locked).}
Deterministic, sub-millisecond, RAM-resident. Holds the rule grammar,
the beam search, the fit/predict pipeline, and the property-binding
resolver. Inference is pure function of input. No I/O during the
predict path. CI-attested Tier-1.

\paragraph{Cortex (SSD-locked).}
Episodic memory, retrieval over prior task signatures, and the
provenance-enforced memory contract that gates write access. Cortex
contributes a seed prior to the Reflex's beam search; under strict
Tier-1, Cortex's contribution is mechanical (no learned weights,
no LLM) and disclosed at the manifest layer.

\paragraph{Contract location.}
The property-bound contract lives in Reflex. Its expressiveness is
the binding axis of the Reflex's competence ceiling. Cortex changes
\emph{which} rules get tried first; the contract changes \emph{which}
rules can fit at all. The bottleneck we report
(\S\ref{sec:bottleneck}) is a Reflex bottleneck. Cortex-side gains
are orthogonal and bounded above by Reflex expressiveness.

% ============================================================
\section{Limitations and Future Work}

\paragraph{The eval gap.}
Under the property-bound grammar at $65$ factories, our Wave-15 full
measurement landed $42/1000$ on the public training split ($4.2\%$) and
$0/120$ on the public evaluation split. The Wave-10 eval diagnostic
records the underlying mechanism: $120/120$ eval tasks had
\emph{zero} rules fitting at all, and consequently
$0/120$ predicted correctly \cite{wave10eval}. The training-to-eval
collapse is not a fit-failure of any particular rule; it is the eval set
falling entirely outside the relations our 65-rule grammar can express.
This is a clean diagnostic: it isolates out-of-distribution
generalization, not within-distribution measurement noise, as the
limiting factor.

\paragraph{Strict Tier-1 ceiling.}
We do not claim that property-binding clears the Tier-1 ceiling on
ARC-AGI-2 eval. We claim it is necessary, not sufficient. Sufficiency
likely requires (i) per-task program synthesis with explicit backtracking
over rule-type instantiations; (ii) richer relational primitives over
multi-hop neighbor structure, periodicity, and partial-pattern
continuation; (iii) an object-identity layer above pixel-level rules so
that ``the subject object'' can be a stable referent across train pairs
even when its concrete properties vary \cite{innovations}.

\paragraph{Hardware roadmap.}
The property-bound resolver is a per-test-input property scan over a
small fixed grid. It is a natural fit for memristor compute-in-memory
substrates and for direct sampling from the test input via a Tsetlin-
Sequencer-Unit (TSU) front-end, both of which are on the Federation's
2027 hardware roadmap. We do not measure them here.

% ============================================================
\section{Conclusion}

The property-bound rule contract is the foundational primitive for
Tier-1 strict ARC rule-template solvers. We name the bottleneck the
contract was missing ($980/1000$ no-fit under global-fit), we
demonstrate that contract expressiveness, not vocabulary size, is the
binding constraint (Waves~5, 6: $+0.00$pp; Waves~7, 8, 9: $+0.30$,
$+0.10$, $+0.30$pp), and we report a $\sim 30\times$ R\&D acceleration
from pre-flight/measurement equivalence under deterministic solvers.
The eval-set collapse to $0/120$ tells us the next bottleneck cleanly:
out-of-distribution generalization, not within-distribution coverage.

% ============================================================
\section*{Acknowledgments}

This work is filed under the AtomEons Federation Charter v1.0. The
co-authorship of Claude Opus 4.7 is recorded per Charter
Article~II~\S2.4 (AI Data Right to Be Cited). We thank the ARC Prize
Foundation for the Paper Track venue and for shipping a benchmark that
rewards strict-Tier-1 honesty over leaderboard theater. No external
funding was received.

% ============================================================
\begin{thebibliography}{9}

\bibitem{Chollet2024}
F.~Chollet, M.~Knoop, G.~Kamradt, B.~Landers.
\newblock {ARC-AGI-2: A New Challenge for Frontier AI Reasoning Systems}.
\newblock ARC Prize Foundation, 2024.

\bibitem{SpelkeKinzler2007}
E.~S.~Spelke, K.~D.~Kinzler.
\newblock Core Knowledge.
\newblock {\em Developmental Science}, 10(1):89--96, 2007.

\bibitem{Greenblatt2024}
R.~Greenblatt.
\newblock Getting 50\% on ARC-AGI with GPT-4o.
\newblock Redwood Research note, 2024.

\bibitem{Hodel2023}
M.~Hodel.
\newblock Domain-Specific Language for the Abstraction and Reasoning Corpus.
\newblock arXiv:2305.18354, 2023.

\bibitem{Icecuber2020}
Icecuber.
\newblock First-place solution to the 2020 ARC Kaggle competition.
\newblock Kaggle write-up, 2020.

% --- Receipts and internal artifacts (cited as evidence, not prior art) ---

\bibitem{innovations}
Atom McCree, Claude Opus 4.7.
\newblock Innovations and Alpha, Waves 4--9.
\newblock Internal artifact,
\texttt{C:\textbackslash AtomEons\textbackslash arc-agi-3-misfit-agent\textbackslash docs\textbackslash INNOVATIONS\_WAVE\_4\_THROUGH\_9.md},
2026-06-16.

\bibitem{wave7receipt}
Atom McCree.
\newblock Wave 7 Orange3 Receipt: property-bound contract, $+0.30$pp lift.
\newblock Receipt artifact,
\texttt{receipts/100day/wave7\_orange3\_receipt.json}, 2026-06-16.

\bibitem{wave10eval}
Atom McCree.
\newblock Wave 10 Eval Diagnostic: $0/120$ under 65-factory grammar.
\newblock Receipt artifact,
\texttt{receipts/100day/wave10\_eval\_diagnostic.json}, 2026-06-17.

\bibitem{blackmamba}
Atom McCree, Claude Opus 4.7.
\newblock Black Mamba v2: Reflex/Cortex Double-Brain Substrate.
\newblock AtomEons Federation Reference Implementation, 2026.

\end{thebibliography}

\end{document}
